\section{Пробный запуск: постановка и порядок проверки}
\label{sec:trial}

\subsection{Назначение локального стенда}

Пробный запуск предназначен для проверки согласованности трёх компонентов:
кодирования входов, выполнения операций ABY3 в TF~Encrypted и
восстановления метки класса. Предусматривается размещение трёх логических
вычислительных сторон на одном компьютере. Такой стенд позволяет сверять
результат после каждого тестового сценария. Независимое администрирование
трёх серверов, предусмотренное моделью нарушителя, требует отдельной
распределённой конфигурации.

Для воспроизведения запуска в протокол испытания включаются сведения
об операционной системе, архитектуре процессора, объёме памяти, версиях
Python, TensorFlow и TF~Encrypted. Перед загрузкой дерева требуется
проверить доступность реализации ABY3, назначение трёх различных ролей
и согласованность типов входных и выходных значений. Масштабирование
фиксируется в одном месте: при передаче уже закодированного целого значения
повторное умножение на масштаб изменило бы операнд сравнения.

\subsection{Контрольные сценарии}

Сначала проверяются операции над разделёнными тестовыми значениями.
Для операндов $-2$ и $3$ ожидаемые результаты сложения и умножения равны
соответственно $1$ и $-6$. Далее проверяется узел с порогом $0$:
при входах $-1$, $0$ и $1$ условие перехода в левую ветвь должно дать
значения $1$, $1$ и $0$. Для каждого сценария восстанавливается только
оговорённый тестовый выход. Эти проверки локализуют ошибки знаковой
интерпретации и обработки равенства до исполнения полного дерева.

Следующий сценарий задаёт дерево глубины два с двумя числовыми
признаками. В корне проверяется условие $x_1\leq 0$. В левом поддереве
проверяется $x_2\leq -1$ с кодами листьев $0$ и $1$, в правом~---
$x_2\leq 1$ с кодами $2$ и $3$. При $f=2$ все указанные пороги и целые
тестовые входы представимы точно. Для пакета
\[
  ((-1,-2),\ (0,0),\ (1,1),\ (1,2))
\]
ожидаемая последовательность меток равна $(0,1,2,3)$. Последовательность
должна сохраняться при поочерёдной подаче тех же объектов. Повторная
обработка пакета со свежими долями должна давать те же открытые метки.

Отдельный пример проверяет влияние округления. Для входа
$(1/16,0)$ открытый обход исходного дерева даёт класс $2$. При $f=2$
первая координата округляется к нулю, поэтому целочисленный эталон
даёт класс $1$. Ожидаемый результат ABY3 здесь равен $1$: его совпадение
с целочисленным эталоном одновременно с расхождением с исходной моделью
проверяет трактовку двух показателей раздела~\ref{sec:methodology}.

\subsection{Регистрация результата}

Для каждого сценария предусматривается запись версии модели, значения
$f$, открытого синтетического входа, ожидаемой и фактически полученной
меток. При расхождении фиксируется первый сценарий, на котором оно
возникло; для диагностики отдельно повторяются сравнение и выбор
дочернего значения на тех же синтетических операндах. Критерием успешного
пробного запуска служит совпадение всех фактических выходов с указанными
эталонами, включая порядок меток в пакете. Приведённые числовые значения
получены аналитически и предназначены для последующей сверки при исполнении
прототипа.
